\section{Diseño de Datos (IEEE 1016)}

Aunque el sistema es \textit{stateless} y no administra bases de datos relacionales tradicionales como PostgreSQL ni no relacionales como MongoDB, existe un modelo de datos lógico muy riguroso para la transferencia de información (Data Transfer Objects - DTOs) vía HTTP. Toda la comunicación entre la capa de negocio y la API LLM se rige por la sintaxis estándar de OpenAI \textit{Chat Completions} adoptada por Groq.

\subsection{Estructuras de Datos de Petición (Payload JSON)}
Cuando el usuario presiona "Generar" (por ejemplo, en el módulo de Objetivos), la vista envía al \texttt{groqService} un objeto estructurado. El servicio ensambla la siguiente estructura JSON sobre el protocolo HTTP POST.

\begin{verbatim}
POST /openai/v1/chat/completions
Content-Type: application/json
Authorization: Bearer VITE_GROQ_API_KEY

{
  "model": "llama3-70b-8192", // Modelo optimizado (70B parámetros)
  "temperature": 0.5,         // Baja entropía, mayor determinismo
  "messages": [
    {
      "role": "system",
      "content": "Eres un tutor académico de tesis..."
    },
    {
      "role": "user",
      "content": "Formula objetivos. Tema: IA en Salud. 
                  Propósito: Mejorar diagnósticos."
    }
  ]
}
\end{verbatim}

\textbf{Justificación de variables técnicas:}
\begin{itemize}
    \item \textbf{Modelo (llama3-70b-8192):} Se eligió Meta LLaMA 3 alojado en Groq debido a sus excelentes capacidades de razonamiento textual y su ventana de contexto de 8192 tokens, suficiente para procesar formulaciones extensas sin truncamiento.
    \item \textbf{Temperature (0.5):} En \textit{Machine Learning}, una temperatura de 0 resultará en el texto más predecible (ideal para código), mientras que un 1 lo hace altamente creativo (ideal para cuentos). Se ajustó a 0.5 para permitir cierta creatividad al proponer sin desviarse a respuestas sin coherencia académica.
\end{itemize}

\subsection{Estructuras de Respuesta y Estado (Respuesta JSON)}
El modelo responde con el siguiente formato, del cual la interfaz debe extraer exclusivamente el nodo de contenido (\texttt{choices[0].message.content}):

\begin{verbatim}
{
  "id": "chatcmpl-92a...",
  "object": "chat.completion",
  "created": 1722881144,
  "model": "llama3-70b-8192",
  "choices": [
    {
      "index": 0,
      "message": {
        "role": "assistant",
        "content": "### Objetivo General...\n\n### Específicos..."
      },
      "finish_reason": "stop"
    }
  ]
}
\end{verbatim}
El frontend capta este string (Markdown) y lo aloja en una variable de estado inmutable (\texttt{const [resultado, setResultado]}). Al cambiar este estado, el React DOM vuelve a renderizar, provocando que la UI dibuje el texto resultante a través del componente \texttt{<ReactMarkdown>}.
